\documentclass[french,12pt,a4paper]{book}
\usepackage{gensymb} 
\usepackage[utf8]{inputenc}
\usepackage[T1]{fontenc}
\usepackage[english]{babel}
\usepackage[default, scale=.95]{opensans} % police Open Sans
\usepackage{amsmath}
\usepackage{amsfonts}
\usepackage{fancyhdr}
\usepackage{amssymb}
\usepackage{xcolor} % où color selon l'installation
\definecolor{Prune}{RGB}{99,0,60} % l14-33 : couleurs de la charte graphique upsaclay
\definecolor{B1}{RGB}{49,62,72} 
\definecolor{C1}{RGB}{124,135,143}
\definecolor{D1}{RGB}{213,218,223}
\definecolor{A2}{RGB}{198,11,70}
\definecolor{B2}{RGB}{237,20,91}
\definecolor{C2}{RGB}{238,52,35}
\definecolor{D2}{RGB}{243,115,32}
\definecolor{A3}{RGB}{124,42,144}
\definecolor{B3}{RGB}{125,106,175}
\definecolor{C3}{RGB}{198,103,29}
\definecolor{D3}{RGB}{254,188,24}
\definecolor{A4}{RGB}{0,78,125}
\definecolor{B4}{RGB}{14,135,201}
\definecolor{C4}{RGB}{0,148,181}
\definecolor{D4}{RGB}{70,195,210}
\definecolor{A5}{RGB}{0,128,122}
\definecolor{B5}{RGB}{64,183,105}
\definecolor{C5}{RGB}{140,198,62}
\definecolor{D5}{RGB}{213,223,61}
\usepackage{mdframed}
\usepackage{multirow} %% Pour mettre un texte sur plusieurs rangées
\usepackage{multicol} %% Pour mettre un texte sur plusieurs colonnes
\usepackage{scrextend} %Forcer la 4ème  de couverture en page pair
\usepackage{tikz}
\usepackage{graphicx}
\usepackage[absolute]{textpos} 
\usepackage{colortbl}
\usepackage{array}
\usepackage{geometry}
\usepackage{titlesec}
\usepackage{tocloft}
\usepackage{hyperref}
\hypersetup{ % paramétrage couleur des liens hypertextes, toujours garder colorlinks=true
	colorlinks=true,
	linkcolor=black,
	urlcolor=purple}
 
\makeatletter
\renewcommand{\numberline}[1]{%
  \@cftbsnum #1\@cftasnum~\@cftasnumb%
}
\makeatother

\pagestyle{plain} % pour ne garder que les n°de page en milieu-bas et supprimer les indications de chapitre en marge haute

\usepackage{lipsum} % à retirer!!!
\begin{document}
	
	\begin{titlepage}
		
		%\thispagestyle{empty}
		
		\newgeometry{left=6cm,bottom=2cm, top=1cm, right=1cm}
		
		\tikz[remember picture,overlay] \node[opacity=1,inner sep=0pt] at (-13mm,-135mm){\includegraphics{Logo/Frame-ups.pdf}};
		
		%*****************************************************
		%******** NUMÉRO D'ORDRE DE LA THÈSE À COMPLÉTER *****
		%******** POUR LE SECOND DÉPOT                   *****
		%*****************************************************
		
		\color{white}
		
		\begin{picture}(0,0)
			\put(-152,-743){\rotatebox{90}{\Large \textsc{THESE DE DOCTORAT}}} \\
			\put(-120,-743){\rotatebox{90}{NNT : 2020UPASA001}}
		\end{picture}
		
		%*****************************************************
		%**  LOGO  ÉTABLISSEMENT PARTENAIRE SI COTUTELLE
		%**  CHANGER L'IMAGE PAR DÉFAUT **
		%*****************************************************
		\vspace{-14mm} % à ajuster en fonction de la hauteur du logo
		%\flushright \includegraphics[scale=1]{logo2.png}
		
		%*****************************************************
		%******************** TITRE **************************
		%*****************************************************
		
		\flushright
		\vspace{30mm} % à régler éventuellement
		\color{Prune}
		\fontfamily{cmss}\fontseries{m}\fontsize{22}{26}\selectfont
		\Huge Device-independent certification: from foundations to applications \\
		
		\normalsize
		\color{black}
		\Large{\textit{Device-independent certification: Fondations et Applications}}\\
		%*****************************************************
		
		\fontfamily{fvs}\fontseries{m}\fontsize{8}{12}\selectfont
		
		\vspace{1.5cm}
		
		\normalsize
		\textbf{Thèse de doctorat de l'université Paris-Saclay} \\
		
		\vspace{6mm}
		
		\small École doctorale n\degree564 : physique en Île-de-France (PIF)\\
		\small Spécialité de doctorat: Physique\\
		\small Graduate School : Physique, Référent : Faculté des sciences d'Orsay \\
		\vspace{6mm}
		
		\footnotesize Thèse préparée à l’Institut de Physique Théorique (Université Paris-Saclay, CNRS, CEA) sous la direction de Nicolas Sangouard, directeur de recherche (CEA). \\
		\vspace{15mm}
		
		\textbf{Thèse soutenue à Paris-Saclay, le JJ mois AAAA, par}\\
		\bigskip
		\Large {\color{Prune} \textbf{Xavier VALCARCE}}
		
		%************************************
		\vspace{\fill} % ALIGNER LE TABLEAU EN BAS DE PAGE
		%************************************
		
		\bigskip
		
		\flushleft
		\small \textbf{Composition du jury}\\
		\vspace{2mm}
		\scriptsize
		\begin{tabular}{|p{7cm}l}
			\arrayrulecolor{Prune}
			\textbf{Prénom Nom} &   Président ou Présidente\\ 
			Titre, Affiliation & \\
			\textbf{Prénom Nom} &  Rapporteur \& Examinateur / trice \\ 
			Titre, Affiliation   &   \\ 
			\textbf{Prénom Nom} &  Rapporteur \& Examinateur / trice \\ 
			Titre, Affiliation  &   \\ 
			\textbf{Prénom Nom} &  Examinateur ou Examinatrice \\ 
			Titre, Affiliation   &   \\ 
			\textbf{Prénom Nom} &  Examinateur ou Examinatrice \\ 
			Titre, Affiliation   &   \\ 
			\textbf{Nicolas Sangouard} &  Directeur ou Directrice de thèse \\ 
			Titre, Affiliation   &   \\ 
			
		\end{tabular} 
		
	\end{titlepage}

\Ifthispageodd{\newpage\thispagestyle{empty}\null\newpage}{}
\thispagestyle{empty}
\newgeometry{top=1.5cm, bottom=1.25cm, left=2cm, right=2cm}
\fontfamily{rm}\selectfont

\lhead{}
\rhead{}
\rfoot{}
\cfoot{}
\lfoot{}

\noindent 
%*****************************************************
%***** LOGO DE L'ED À CHANGER IMPÉRATIVEMENT *********
%*****************************************************
\includegraphics[height=2.45cm]{Logo/logo_usp_PIF.png}
\vspace{1cm}
%*****************************************************
\fontfamily{cmss}\fontseries{m}\selectfont

\small

\begin{mdframed}[linecolor=Prune,linewidth=1]
	
	\textbf{Titre:} titre (en français)..................................................................................................................
	
	\noindent \textbf{Mots clés:} 3 à 6 mots clefs (version en français)
	
	\vspace{-.5cm}
	\begin{multicols}{2}
		\noindent \textbf{Résumé:}\lipsum[1-2] 
	\end{multicols}
	
\end{mdframed}

\vspace{8mm}

\begin{mdframed}[linecolor=Prune,linewidth=1]
	
	\textbf{Title:} titre (en anglais)..................................................................................................................
	
	\noindent \textbf{Keywords:} 3 à 6 mots clefs (version en anglais)
	
	\begin{multicols}{2}
		\noindent \textbf{Abstract:} \lipsum[1-2]
	\end{multicols}
\end{mdframed}

\titleformat{\part}[block]{\bfseries\Huge\color{Prune}\filright}{\thepart\ --\ }{.1ex}
{\vspace{0.1ex}
}
[\vspace{1ex}]

\titleformat{\chapter}[hang]{\bfseries\Large\color{Prune}}{\thechapter\ -}{.1ex}
{\vspace{0.1ex}
}
[\vspace{1ex}]
\titlespacing{\chapter}{0pc}{0ex}{0.5pc}

\titleformat{\section}[hang]{\bfseries\normalsize}{\thesection\ .}{0.5pt}
{\vspace{0.1ex}
}
[\vspace{0.1ex}]
\titlespacing{\section}{1.5pc}{4ex plus .1ex minus .2ex}{.8pc}

\titleformat{\subsection}[hang]{\bfseries\small}{\thesubsection\ .}{1pt}
{\vspace{0.1ex}
}
[\vspace{0.1ex}]
\titlespacing{\subsection}{3pc}{2ex plus .1ex minus .2ex}{.1pc}

\newcounter{article}
\newcommand{\article}[1]{%
  \par\refstepcounter{article}% Increase section counter
  \sectionmark{#1}% Add section mark (header)
  \addcontentsline{toc}{chapter}{\protect\numberline{Article \thearticle}#1}% Add section to ToC
  % Add more content here, if needed.
}
\frontmatter

\newgeometry{top=4cm, bottom=4cm, left=2cm, right=2cm}

\chapter{Acknowledgements}
\newpage

\tableofcontents

\newgeometry{top=4cm, bottom=4cm, left=4cm, right=4cm}

\chapter{Acronyms}

\chapter{Thesis overview}

\section{Motivation and background}

\section{Summary of the contributions}

\chapter{Synthèse en français}

\section{Motivations et context}

\section{Résumé des contributions}

\chapter{Organisation of the manuscript}

\chapter{Publications}

\mainmatter

\part{Self-testing for device-independent certification}

\chapter{Overview and definitions}

\chapter{Self-testing of bipartite states}

\chapter{Robust self-testing}


\article{What is the minimum CHSH score certifying that a state resembles the singlet?}
\article{Self-testing two-qubit maximally entangled states from generalized Clauser-Horne-Shimony-Holt tests}

\part{Device-independent communication}

\end{document}